\section{Short Troubleshooting Order}

When execution fails, check the pipeline in the same order in which it was
built:

\begin{table}[H]
  \centering
  \begin{tabularx}{\textwidth}{@{}p{0.27\textwidth}X@{}}
    \toprule
    \textbf{Symptom} & \textbf{First check} \\
    \midrule
    Scene or reset is wrong & Open the XML alone; inspect initial state, contacts, and randomization. \\
    Teleoperation is delayed & Check Touch topics, newest-target handling, scaling, smoothing, and simulation rate. \\
    Dataset looks wrong & Replay the affected episode; inspect frozen images, episode boundaries, action jumps, and gripper labels. \\
    Checkpoint will not load & Match source version, task YAML, policy class, and Python dependencies. \\
    Policy moves to the wrong pose & Check image order, normalization, action frame, quaternion convention, and reset distribution. \\
    Real robot tracks poorly & Compare raw target with measured pose, then inspect latency, controller limits, and safety output. \\
    \bottomrule
  \end{tabularx}
\end{table}

Change one variable at a time and repeat the smallest test that can confirm the
cause. Do not collect a larger dataset until scene, teleoperation, recording,
and deployment interfaces have been checked.

\section{What to Leave for the Next Student}

Keep a short record containing the dataset path and episode count, matching
training configuration, selected checkpoint, exact execution command, observed
inference latency, and known failure cases. With those items and the three
repository README files, another student can repeat the workflow without
reconstructing it from terminal history.

\begin{checkpointbox}[Complete workflow]
The task is stable in MuJoCo, teleoperation is responsive, accepted episodes
have been replayed, the checkpoint matches the dataset interface, and the policy
has been executed successfully in simulation before any physical deployment.
\end{checkpointbox}
